\documentclass[11pt,a4paper]{article}

% Packages
\usepackage[utf8]{inputenc}
\usepackage[T1]{fontenc}
\usepackage[margin=0.75in]{geometry}
\usepackage{enumitem}
\usepackage{hyperref}
\usepackage{titlesec}
\usepackage{xcolor}

% Formatting
\pagestyle{empty}
\setlength{\parindent}{0pt}
\setlist[itemize]{leftmargin=*, noitemsep, topsep=3pt}

% Section formatting
\titleformat{\section}{\Large\bfseries}{}{0em}{}[\titlerule]
\titlespacing*{\section}{0pt}{12pt}{6pt}

\titleformat{\subsection}[runin]{\bfseries}{}{0em}{}
\titlespacing*{\subsection}{0pt}{8pt}{0pt}

% Colors
\definecolor{linkcolor}{RGB}{0,102,204}
\hypersetup{
    colorlinks=true,
    urlcolor=linkcolor
}

\begin{document}

% Header
\begin{center}
    {\LARGE \textbf{ERIC MAYEUX}}\\
    \vspace{4pt}
    {\large Gestionnaire PMO -- NPI \& Gouvernance R\&D | PMO Manager NPI \& R\&D Governance}\\
    \vspace{4pt}
    Montréal, QC, Canada\\
    \href{mailto:mayeux.e@gmail.com}{mayeux.e@gmail.com} | 438-884-7645
\end{center}

\vspace{8pt}

\section*{Profil Professionnel}
\noindent
\textbf{Gestionnaire PMO} avec 10+ ans d'expérience en \textbf{gouvernance de projets}, standardisation des processus NPI et \textbf{gestion de portefeuille R\&D} en environnement technique. Expert en définition de cadres de gouvernance, \textbf{indicateurs de performance (KPI)}, tableaux de bord et \textbf{rapports structurés} pour la direction. Maîtrise de Jira, MS Project et méthodologies \textbf{Agile/Scrum}. Certifié \textbf{PMP} (en cours), CSM. Bilinguisme français/anglais.

\section*{Compétences Techniques}

\textbf{Gouvernance PMO \& NPI:} Normes de gestion de projets, cadres de gouvernance, jalons, KPI, standardisation NPI, cycle de vie produit, amélioration continue du cadre PMO

\textbf{Gestion de portefeuille:} Vue globale du portefeuille R\&D, priorisation, allocation des ressources, dépendances, rapports structurés à la direction

\textbf{Outils PMO:} Jira, MS Project, PLM, tableaux de bord, Confluence, Datadog, Splunk

\textbf{Analyse \& Reporting:} Synthèse d'informations complexes, rapports à la direction, traçabilité des données, suivi financier (CapEx/OpEx)

\textbf{Méthodologies:} Agile, Scrum, SAFe, Kanban, Waterfall, PMP, gestion du cycle de vie produit

\textbf{Intelligence Artificielle:} Agents IA, automatisation, RPA, Playwright MCP, Crew AI

\textbf{Architecture \& Cloud:} AWS Cloud Practitioner, CI/CD, Big Data

\vspace{6pt}

\section*{Réalisations}
\begin{itemize}
    \item Définition et mise en œuvre de \textbf{normes et cadres de gouvernance} pour le portefeuille de projets TI
    \item Mise en place de \textbf{jalons, KPI et standards de rapports} pour la direction et les parties prenantes
    \item \textbf{Standardisation des processus} de livraison et amélioration de la prévisibilité des projets
    \item Consolidation d'une \textbf{vue globale du portefeuille} : santé de livraison, capacité, dépendances
    \item Déploiement d'outils de \textbf{gestion de portefeuille} (Jira, tableaux de bord) et leçons apprises
    \item \textbf{Automatisation} des tâches administratives et de reporting grâce à des agents IA
\end{itemize}

\section*{Expérience Professionnelle}

\subsection*{Scrum Master -- Gestion de projet | Project Manager} \hfill \textit{Novembre 2025 -- Présent}\\
\textbf{Banque Nationale du Canada (BNC)} -- Montréal, QC\\
Définit et améliore les \textbf{normes et cadres de gouvernance} du portefeuille de projets TI.
\begin{itemize}
    \item Établit et maintient la \textbf{gouvernance des projets} : jalons, KPI, normes de rapports
    \item Standardise les \textbf{processus de livraison} pour toutes les équipes du train (SAFe)
    \item Consolide une \textbf{vue globale du portefeuille} : santé de livraison, capacité, dépendances, finances
    \item Fournit des \textbf{rapports structurés} à la direction ; soutient la priorisation et l'allocation des ressources
    \item Intègre les leçons apprises et l'\textbf{amélioration continue} dans les pratiques de livraison
    \item Déploie des \textbf{agents IA} pour réduire la charge administrative des équipes
\end{itemize}

\subsection*{Intégrateur Sénior Qualité | Senior Quality Integrator} \hfill \textit{Octobre 2023 -- Novembre 2025}\\
\textbf{Banque Nationale du Canada (BNC)} -- Montréal, QC\\
\begin{itemize}
    \item Définit les \textbf{normes PMO} qualité : métriques, jalons, traçabilité des données (Datadog, Splunk)
    \item Gère le \textbf{portefeuille} des initiatives qualité pan-train ; allocation des ressources
    \item \textbf{Rapports structurés} à la haute direction ; coaching et formation continue
\end{itemize}

\subsection*{Intégrateur Qualité -- SDET | Quality Integrator -- Software Development Engineer in Test} \hfill \textit{Janvier 2022 -- Octobre 2023}\\
\textbf{Banque Nationale du Canada (BNC)} via Groupe Astek -- Montréal, QC\\
\begin{itemize}
    \item Mise en place de métriques qualité et \textbf{cadre PMO} : pyramide de test, couverture, sécurité
    \item Stack : Selenium, AppliTools, RestAssured, JMeter, Gatling, K6 ; tableaux de bord Splunk, Jira
\end{itemize}

\subsection*{QA Automaticien | QA Automation Engineer} \hfill \textit{Juin 2021 -- Décembre 2021}\\
\textbf{AODocs} -- Paris, France\\
\begin{itemize}
    \item Automatisation tests ; POC Flutter ; rapports Allure
\end{itemize}

\subsection*{Automaticien CI/CD | DevOps Automation Engineer} \hfill \textit{Janvier 2020 -- Juin 2021}\\
\textbf{ENGIE DIGITAL} via Groupe Hénix -- Paris, France\\
\begin{itemize}
    \item Création CI/CD ; \textbf{administration Jira}, KPIs, \textbf{gouvernance projets} ; scripting Python/Shell
\end{itemize}

\subsection*{Product Owner} \hfill \textit{Juin 2019 -- Septembre 2019}\\
\textbf{Hénix} -- Montrouge, France\\
\begin{itemize}
    \item Gestion du backlog, priorisation, certification Jira ACP-600 (8 500 utilisateurs)
\end{itemize}

\subsection*{Développeur Logiciel | Software Developer} \hfill \textit{Mai 2017 -- Mai 2019}\\
\textbf{MEDIAPOST} via Groupe Hénix -- Paris, France\\
\begin{itemize}
    \item Développement back office, webservices, scripts BDD ; reporting et macros Excel
\end{itemize}

\subsection*{Consultant QA \& Automatisation | QA Automation Consultant} \hfill \textit{Avril 2016 -- Avril 2017}\\
\textbf{ENGIE IT} via Groupe Hénix -- Saint-Ouen, France\\
\begin{itemize}
    \item Gestion opérationnelle des environnements, coordination fournisseurs, administration HP ALM
\end{itemize}

\section*{Certifications \& Formations}

\textbf{PMP Certification} (en cours) \hfill \textit{2026}

Project Management Professional -- Project Management Institute (PMI)

\textbf{AWS Cloud Practitioner} \hfill \textit{2026}

Amazon Web Services (AWS)

\textbf{Certified Scrum Master (CSM)} \hfill \textit{2024}

Scrum Alliance

\textbf{Jira ACP-600 -- Project Administration in Jira Server} \hfill \textit{2019}

Atlassian Certified Professional

\textbf{Cursus Automaticien et DevOps} \hfill \textit{2019}\\
\textit{AFCEPF -- Montrouge, France}

\textbf{Cursus Architecture logiciel \& Analyste informaticien - Certifié Master II} \hfill \textit{2017}\\
\textit{AFCEPF-HENIX -- Malakoff, France}

\section*{Formation Académique}

\textbf{Licence en Gestion (Bachelor's Degree in Management)} \hfill \textit{2011}\\
École Supérieure de Gestion (E.S.G) -- Paris, France

\section*{Compétences Linguistiques}

\textbf{Français:} Langue maternelle (Native)\\
\textbf{Anglais:} Niveau Avancé (Advanced / Professional Working Proficiency)

\end{document}
